\chapter{Coordinate Geometry}\label{ch:g10:coordgeom}

Descartes' great idea was to describe points by pairs of numbers, turning
geometry problems into computations. With just two formulas --- the
\omterm{prop:g10:coordgeom:midpoint}{midpoint} and the distance --- one can prove that a triangle is isosceles,
that a quadrilateral is a parallelogram, or that three points are aligned,
without drawing a single auxiliary line.

\section{Coordinates in the plane}

\begin{definition}[Coordinate system]\label{def:g10:coordgeom:system}
A \emph{coordinate system}\index{coordinate system} of the plane consists
of an origin $O$ and two graduated axes through $O$: the horizontal
$x$-axis and the vertical $y$-axis. Every point $M$ then has a unique pair
of \emph{coordinates}\index{coordinates} $(x, y)$: its
\emph{abscissa}\index{abscissa} $x$ and its
\emph{ordinate}\index{ordinate} $y$. The system is
\emph{orthonormal}\index{orthonormal system} when the axes are
perpendicular and carry the same unit of length. All systems in this
chapter are orthonormal.
\end{definition}

\begin{omfigure}
\begin{tikzpicture}[scale=0.85]
  \draw[thin, gray!40] (-2.4,-1.4) grid (4.4,3.4);
  \draw[-Stealth] (-2.6,0) -- (4.6,0) node[below, font=\small] {$x$};
  \draw[-Stealth] (0,-1.6) -- (0,3.6) node[left, font=\small] {$y$};
  \node[below left, font=\small] at (0,0) {$O$};
  \fill[omDef] (3,2) circle (2pt);
  \node[above right, font=\small, omDef] at (3,2) {$M(3, 2)$};
  \draw[dashed, omDef] (3,0) -- (3,2) -- (0,2);
  \fill[omThm] (-2,-1) circle (2pt);
  \node[below left, font=\small, omThm] at (-2,-1) {$N(-2, -1)$};
  \foreach \x in {1,2,3,4} \node[below, font=\footnotesize] at (\x,0) {\x};
  \foreach \y in {1,2,3} \node[left, font=\footnotesize] at (0,\y) {\y};
\end{tikzpicture}

{\small Every point of the plane is located by two numbers: first the
\omterm{def:g10:coordgeom:system}{abscissa} (horizontal), then the \omterm{def:g10:coordgeom:system}{ordinate} (vertical).}
\end{omfigure}

\section{Midpoint of a segment}

\begin{proposition}[Midpoint formula]\label{prop:g10:coordgeom:midpoint}
Let $A(x_A, y_A)$ and $B(x_B, y_B)$. The \emph{midpoint}\index{midpoint}
$I$ of the segment $[AB]$ has \omterm{def:g10:coordgeom:system}{coordinates}
\[
I\left(\frac{x_A + x_B}{2},\ \frac{y_A + y_B}{2}\right).
\]
\end{proposition}

\begin{proof}
Consider the horizontal \omterm{def:g10:coordgeom:system}{coordinates}. Going from $A$ to $B$, the \omterm{def:g10:coordgeom:system}{abscissa}
changes by $x_B - x_A$; the point halfway has \omterm{def:g10:coordgeom:system}{abscissa}
\[
x_A + \frac{x_B - x_A}{2} = \frac{2x_A + x_B - x_A}{2}
= \frac{x_A + x_B}{2},
\]
the average of the two \omterm{def:g10:coordgeom:system}{abscissas}. The same computation applies to the
\omterm{def:g10:coordgeom:system}{ordinates}.
\end{proof}

\begin{example}\label{ex:g10:coordgeom:midpoint}
With $A(-1, 4)$ and $B(5, -2)$, the \omterm{prop:g10:coordgeom:midpoint}{midpoint} of $[AB]$ is
\[
I\left(\frac{-1 + 5}{2},\ \frac{4 + (-2)}{2}\right) = I(2, 1).
\]
The formula also runs backwards: if $A(-1, 4)$ and the \omterm{prop:g10:coordgeom:midpoint}{midpoint} is
$I(2, 1)$, then $B$ satisfies $\frac{-1 + x_B}{2} = 2$ and
$\frac{4 + y_B}{2} = 1$, so $x_B = 5$ and $y_B = -2$: the point $B$ is the
\emph{symmetric} of $A$ about $I$.
\end{example}

\section{Distance between two points}

\begin{theorem}[Distance formula]\label{thm:g10:coordgeom:distance}
In an \omterm{def:g10:coordgeom:system}{orthonormal system}, the distance between $A(x_A, y_A)$ and
$B(x_B, y_B)$ is
\[
AB = \sqrt{(x_B - x_A)^2 + (y_B - y_A)^2}.
\]
\end{theorem}

\begin{proof}
Let $C$ be the point $(x_B, y_A)$: it has the same height as $A$ and the
same \omterm{def:g10:coordgeom:system}{abscissa} as $B$, so the triangle $ACB$ has a right angle at $C$. Its
legs are horizontal and vertical segments, of lengths
$AC = \abs{x_B - x_A}$ and $CB = \abs{y_B - y_A}$. By the Pythagorean
theorem,
\[
AB^2 = AC^2 + CB^2 = (x_B - x_A)^2 + (y_B - y_A)^2,
\]
and taking the square root (both sides are nonnegative) gives the formula.
\end{proof}

\begin{omfigure}
\begin{tikzpicture}[scale=0.85]
  \draw[thin, gray!40] (-0.4,-0.4) grid (5.4,3.4);
  \draw[-Stealth] (-0.6,0) -- (5.6,0) node[below, font=\small] {$x$};
  \draw[-Stealth] (0,-0.6) -- (0,3.6) node[left, font=\small] {$y$};
  \coordinate (A) at (1,1);
  \coordinate (B) at (5,3);
  \coordinate (C) at (5,1);
  \draw[omDef, thick] (A) -- (B);
  \draw[omThm, thick] (A) -- (C) node[midway, below, font=\small]
    {$\abs{x_B - x_A}$};
  \draw[omThm, thick] (C) -- (B) node[midway, right, font=\small]
    {$\abs{y_B - y_A}$};
  \draw[thin] (C) ++(-0.25,0) -- ++(0,0.25) -- ++(0.25,0);
  \fill (A) circle (2pt) node[above left, font=\small] {$A$};
  \fill (B) circle (2pt) node[above, font=\small] {$B$};
  \fill (C) circle (2pt) node[below right, font=\small] {$C$};
\end{tikzpicture}

{\small The distance formula is the Pythagorean theorem applied to the
right triangle $ACB$ built on a horizontal and a vertical leg.}
\end{omfigure}

\begin{example}\label{ex:g10:coordgeom:distance}
With $A(1, 1)$ and $B(5, 3)$:
\[
AB = \sqrt{(5-1)^2 + (3-1)^2} = \sqrt{16 + 4} = \sqrt{20} = 2\sqrt 5 .
\]
Distances are usually kept in exact (square root) form; only round at the
very end if a decimal answer is needed.
\end{example}

\begin{remark}
The order of the points does not matter: $(x_A - x_B)^2 =
(x_B - x_A)^2$. But do not forget the squares! The distance is
\emph{not} $\abs{x_B - x_A} + \abs{y_B - y_A}$.
\end{remark}

\section{Using the formulas in geometry}

\begin{method}[Nature of a triangle]\label{met:g10:coordgeom:triangle}
Given three points $A$, $B$, $C$ by their \omterm{def:g10:coordgeom:system}{coordinates}:
\begin{enumerate}
  \item compute the three squared lengths $AB^2$, $AC^2$, $BC^2$ with the
        distance formula (keeping squares avoids square roots);
  \item two equal squared lengths $\Rightarrow$ the triangle is
        \emph{isosceles};
  \item if the largest squared length is the sum of the two others, the
        triangle is \emph{right-angled} at the vertex opposite the longest
        side, by the converse of the Pythagorean theorem.
\end{enumerate}
\end{method}

\begin{example}\label{ex:g10:coordgeom:triangle}
Let $A(0, 1)$, $B(4, 3)$ and $C(2, -3)$. Then
\begin{align*}
AB^2 &= (4-0)^2 + (3-1)^2 = 16 + 4 = 20, \\
AC^2 &= (2-0)^2 + (-3-1)^2 = 4 + 16 = 20, \\
BC^2 &= (2-4)^2 + (-3-3)^2 = 4 + 36 = 40 .
\end{align*}
First, $AB^2 = AC^2$, so $AB = AC$: the triangle is isosceles at $A$.
Moreover $AB^2 + AC^2 = 40 = BC^2$, so by the converse of the Pythagorean
theorem the triangle is also right-angled at $A$.
\end{example}

\begin{method}[Recognizing a parallelogram]\label{met:g10:coordgeom:parallelogram}
A quadrilateral $ABCD$ is a parallelogram exactly when its diagonals
$[AC]$ and $[BD]$ have the same \omterm{prop:g10:coordgeom:midpoint}{midpoint}. So: compute both \omterm{prop:g10:coordgeom:midpoint}{midpoints} with
the \omterm{prop:g10:coordgeom:midpoint}{midpoint} formula and compare them.
\end{method}

\begin{example}\label{ex:g10:coordgeom:parallelogram}
Let $A(-1, 0)$, $B(2, 2)$, $C(5, 1)$ and $D(2, -1)$. The \omterm{prop:g10:coordgeom:midpoint}{midpoint} of
$[AC]$ is $\left(\frac{-1+5}{2}, \frac{0+1}{2}\right) = (2, \frac12)$; the
\omterm{prop:g10:coordgeom:midpoint}{midpoint} of $[BD]$ is $\left(\frac{2+2}{2}, \frac{2-1}{2}\right) =
(2, \frac12)$. They coincide, so $ABCD$ is a parallelogram.
\end{example}

\begin{omfigure}
\begin{tikzpicture}[scale=0.75]
  \draw[thin, gray!40] (-1.4,-1.4) grid (5.4,2.4);
  \draw[-Stealth] (-1.6,0) -- (5.6,0);
  \draw[-Stealth] (0,-1.6) -- (0,2.6);
  \coordinate (A) at (-1,0);
  \coordinate (B) at (2,2);
  \coordinate (C) at (5,1);
  \coordinate (D) at (2,-1);
  \draw[omDef, thick] (A) -- (B) -- (C) -- (D) -- cycle;
  \draw[omThm, dashed] (A) -- (C);
  \draw[omThm, dashed] (B) -- (D);
  \fill (A) circle (2pt) node[left, font=\small] {$A$};
  \fill (B) circle (2pt) node[above, font=\small] {$B$};
  \fill (C) circle (2pt) node[right, font=\small] {$C$};
  \fill (D) circle (2pt) node[below, font=\small] {$D$};
  \fill[omThm] (2,0.5) circle (2pt) node[above right, font=\small] {$I$};
\end{tikzpicture}

{\small $ABCD$ is a parallelogram because its two diagonals share the same
\omterm{prop:g10:coordgeom:midpoint}{midpoint} $I(2, \frac12)$.}
\end{omfigure}

\section{Exercises}

\begin{exercise}[$\star$]\label{exo:g10:coordgeom:1}
Let $A(2, 5)$ and $B(-4, 1)$. Compute the \omterm{def:g10:coordgeom:system}{coordinates} of the \omterm{prop:g10:coordgeom:midpoint}{midpoint} of
$[AB]$ and the distance $AB$.
\end{exercise}

\begin{exercise}[$\star$]\label{exo:g10:coordgeom:2}
Let $A(3, -2)$ and $I(1, 2)$. Find the \omterm{def:g10:coordgeom:system}{coordinates} of the point $B$ such
that $I$ is the \omterm{prop:g10:coordgeom:midpoint}{midpoint} of $[AB]$.
\end{exercise}

\begin{exercise}[$\star$]\label{exo:g10:coordgeom:3}
Which of the points $P(4, 1)$, $Q(-3, 2)$ and $R(0, -5)$ is closest to the
origin $O(0,0)$? Answer with exact distances.
\end{exercise}

\begin{exercise}[$\star$]\label{exo:g10:coordgeom:4}
Plot the points $A(1, 2)$, $B(5, 4)$, $C(7, 0)$ on \omterm{def:g10:functions:graph}{graph} paper, then show
by computation that the triangle $ABC$ is isosceles. Is it right-angled?
\end{exercise}

\begin{exercise}[$\star\star$]\label{exo:g10:coordgeom:5}
Let $A(-2, 1)$, $B(1, 3)$, $C(4, 1)$ and $D(1, -1)$.
\begin{enumerate}
  \item Show that $ABCD$ is a parallelogram.
  \item Compute $AB$ and $BC$. Is $ABCD$ a rhombus (all sides equal)?
  \item Compute $AC$ and $BD$. Is $ABCD$ a rectangle (equal diagonals)?
\end{enumerate}
\end{exercise}

\begin{exercise}[$\star\star$]\label{exo:g10:coordgeom:6}
The circle with center $\Omega(2, 1)$ and radius $5$ consists of all
points at distance $5$ from $\Omega$. Which of the points $A(5, 5)$,
$B(-1, -3)$, $C(6, 2)$ lie on this circle? Inside it? Outside it?
\end{exercise}

\begin{exercise}[$\star\star$]\label{exo:g10:coordgeom:7}
Let $A(1, 1)$ and $B(7, 5)$. Find all points $M(x, 0)$ of the $x$-axis
that are equidistant from $A$ and $B$. (Write $MA^2 = MB^2$ and solve for
$x$.)
\end{exercise}

\begin{exercise}[$\star\star$]\label{exo:g10:coordgeom:8}
Let $A(0, 3)$, $B(4, 1)$.
\begin{enumerate}
  \item Compute the \omterm{def:g10:coordgeom:system}{coordinates} of the \omterm{prop:g10:coordgeom:midpoint}{midpoint} $I$ of $[AB]$.
  \item Show that $M(2, 2) = I$, then verify by computing $MA$, $MB$ that
        $M$ is equidistant from $A$ and $B$.
  \item Find a second point, on the $y$-axis, equidistant from $A$ and
        $B$.
\end{enumerate}
\end{exercise}

\begin{exercise}[$\star\star$]\label{exo:g10:coordgeom:9}
The points $A(-1, -1)$, $B(3, 1)$ and $C(11, 5)$ are given. Compute $AB$,
$BC$ and $AC$, and deduce that $A$, $B$, $C$ are aligned. (Three points
are aligned exactly when the largest of the three distances is the sum of
the two others.)
\end{exercise}

\begin{exercise}[$\star\star\star$]\label{exo:g10:coordgeom:10}
Let $A(a, 0)$ and $B(0, b)$ with $a, b > 0$, and let $I$ be the \omterm{prop:g10:coordgeom:midpoint}{midpoint}
of $[AB]$. Show by computation that $OI = IA = IB$, where $O$ is the
origin. (This proves a classical theorem: in a right triangle, the
\omterm{prop:g10:coordgeom:midpoint}{midpoint} of the hypotenuse is equidistant from the three vertices.)
\end{exercise}

\section{Problem: Descartes' bridge --- circle equations and how
to find where you are}

\begin{problem}[{Weekend problem --- the circle becomes an
equation, old theorems become computations, and three beacons
locate a point: trilateration by algebra}]
\label{pb:g10:coordgeom:1}
In 1637 Ren\'e Descartes built a bridge between two continents:
every curve of geometry became an \emph{\omterm{def:g10:algebra:equation}{equation}}, every
geometric question a computation. This problem walks the bridge
in both directions --- deriving the circle's \omterm{def:g10:algebra:equation}{equation} from the
distance formula (\cref{thm:g10:coordgeom:distance}), re-proving
classical theorems in three lines of algebra, and ending where
the bridge carries the most traffic today: computing a position
from distance signals, the flat-earth heart of GPS.

\medskip
\noindent\textbf{Part I --- The circle gets an \omterm{def:g10:algebra:equation}{equation}.}
\begin{enumerate}
  \item A point $M(x, y)$ lies on the circle of center
        $\Omega(2, 1)$ and radius $5$ exactly when
        $\Omega M = 5$. Square this condition to obtain the
        circle's \omterm{def:g10:algebra:equation}{equation}:
        \[
        (x - 2)^2 + (y - 1)^2 = 25 .
        \]
  \item Test membership: which of $A(5, 5)$, $B(6, 4)$,
        $D(4, 4)$ are on the circle, inside it, outside it?
  \item A circle in disguise: complete the squares
        (\cref{pb:g10:algebra:1}) in
        \[
        x^2 + y^2 - 6x + 4y - 12 = 0
        \]
        and give its center and radius.
  \item Do the same with $x^2 + y^2 - 6x + 4y + 14 = 0$. What
        does the completed form reveal? State the criterion: when
        does $(x - h)^2 + (y - k)^2 = c$ describe a circle, a
        single point, or nothing at all?
  \item Intersect the circle of question 3 with the horizontal
        line $y = 3$: substitute and solve. How many points, and
        what is the geometric name of such a line?
\end{enumerate}

\medskip
\noindent\textbf{Part II --- Old theorems in three lines.}
\begin{enumerate}[resume]
  \item The circle theorem of the Middle School volume's circle theorem (a right angle in every half-circle),
        re-proved by algebra: let $A(-r, 0)$ and $B(r, 0)$ be
        the ends of a diameter and $M(x, y)$ any point of the
        circle $x^2 + y^2 = r^2$. Compute $MA^2 + MB^2$ and
        conclude with the converse of Pythagoras that
        $\widehat{AMB}$ is right.
  \item The median theorem: with $A(-a, 0)$, $B(a, 0)$ (so the
        \omterm{prop:g10:coordgeom:midpoint}{midpoint} of $[AB]$ is the origin $I$), show that for
        every point $M(x, y)$:
        \[
        MA^2 + MB^2 = 2\,MI^2 + \frac{AB^2}{2} .
        \]
  \item Deduce the locus of the points $M$ with
        $MA^2 + MB^2 = 26$ when $AB = 6$: what curve, which
        center, what radius?
  \item A different locus: $A(0,0)$, $B(3,0)$; find all points
        with $MA = 2\,MB$. (Square, expand, complete the
        squares: a famous circle appears --- Apollonius knew it
        without \omterm{def:g10:coordgeom:system}{coordinates}.)
  \item In one or two sentences: what does Descartes' bridge
        change about \emph{how} theorems get proved? Compare
        question 6 with the rectangle proof of
        the Middle School volume's circle theorem (a right angle in every half-circle).
\end{enumerate}

\medskip
\noindent\textbf{Part III --- Three beacons find you.}
Your position $M(x, y)$ is unknown. Beacon $P(0, 0)$ measures
your distance as $5$; beacon $Q(6, 0)$ also measures $5$.
\begin{enumerate}[resume]
  \item Write the two circle \omterm{def:g10:algebra:equation}{equations}, subtract them, and watch
        the squares cancel: what simple \omterm{def:g10:algebra:equation}{equation} survives?
        Combine with one circle to find the \emph{two} candidate
        positions.
  \item A third beacon $R(0, 8)$ measures your distance as $5$.
        Compute its distance to each candidate and decide where
        you are.
  \item Explain why subtracting two circle \omterm{def:g10:algebra:equation}{equations}
        \emph{always} yields the \omterm{def:g10:algebra:equation}{equation} of a line (which terms
        cancel?), and what that line is geometrically when the
        circles cross at two points.
  \item Real satellite positioning works in space and with
        imperfect clocks: each satellite gives (via signal travel
        time) one distance, hence one sphere. How many unknowns
        does a receiver have (position \emph{and} its own clock
        error), and why does GPS therefore listen to at least
        \emph{four} satellites?
  \item Precision: suppose beacon $P$'s distance is really
        $5.1$ instead of $5$ (all else unchanged). Redo the
        subtraction of question 11 to find the new $x$, and
        quantify how far the estimate moved. (Compare the error
        budget ideas of \cref{pb:g10:numbers:1}.)
\end{enumerate}

\medskip
\noindent\textbf{Part IV --- The surveyor's toolkit.}
\begin{enumerate}[resume]
  \item The shortest watering path: a camp at $A(1, 5)$, a barn
        at $B(7, 3)$, a straight river along the axis $y = 0$.
        To go from $A$ to the river and then to $B$ with the
        least walking: reflect $B$ across the river into $B'$,
        and find where the segment $[AB']$ crosses the river.
        Give the crossing point and the minimal total length.
        (The idea is the billiard-bounce exercise of the Middle School volume, now fully
        computable.)
  \item Confirm minimality on a rival: compute the total path
        via $Q(3, 0)$ and check it loses to your answer.
  \item Classify the quadrilateral $A(1,1)$, $B(4,2)$,
        $C(5,5)$, $D(2,4)$ completely: parallelogram? rhombus?
        rectangle? square?
        (\cref{met:g10:coordgeom:parallelogram},
        \cref{met:g10:coordgeom:triangle} --- compare diagonals
        too.)
  \item Given $A(-1, 2)$, $B(3, 4)$, $C(6, 0)$, find $D$ such
        that $ABCD$ is a parallelogram (the \omterm{prop:g10:coordgeom:midpoint}{diagonal-midpoint}
        trick of the half-turns weekend problem of the Middle School volume, now in formulas).
  \item Finale: three tools --- \omterm{prop:g10:coordgeom:midpoint}{midpoint} formula, distance
        formula, subtraction of circle \omterm{def:g10:algebra:equation}{equations}. For each, name
        the kind of question it settled in this problem, and
        state what the next two chapters add to the kit
        (directions and \omterm{def:g10:reffunc:affine}{slopes}: vectors and line \omterm{def:g10:algebra:equation}{equations}).
\end{enumerate}
\end{problem}
